\section{二项式与双重计数}
\label{sec:04-Discrete-Mathematics/01-Combinatorics/02}

在七个候选人里选三个当班委，你知道答案是 \(C_7^3=35\)。这个数是怎么和 \((x+y)^7\) 的展开联系在一起的？可能你觉得``二项式定理就是一个展开公式''，对着 Pascal 三角查系数就好。那我们来做一个简单的实验。

写下 \((x+y)^3\) 的展开，不套公式，老老实实乘开：

\[
(x+y)^3=(x+y)(x+y)(x+y).
\]

先全部展开，再合并同类项。三个括号，每个括号里可以选 \(x\) 或者选 \(y\)。如果三个括号都选 \(x\)，得到一项 \(x^3\)——只有一种选法，系数是 1。如果恰有一个括号选 \(y\)，其他两个选 \(x\)，得到 \(x^2y\)——从三个括号里挑一个贡献 \(y\)，有 \(C_3^1=3\) 种挑法，所以合并后有 \(3x^2y\)。两个 \(y\) 一个 \(x\)，同理是 \(C_3^2=3\) 种，得到 \(3xy^2\)。三个都选 \(y\)，只有一种，得到 \(y^3\)。

系数 \(1,3,3,1\) 正好是 \(C_3^0,C_3^1,C_3^2,C_3^3\)。这不是巧合。

\subsection{展开式的系数在数``选哪些因子''}

把 \((x+y)^n\) 写成 \(n\) 个因子的乘积：

\[
(x+y)^n=\underbrace{(x+y)(x+y)\cdots(x+y)}_{n\text{ 个因子}}.
\]

要得到一项 \(x^{n-k}y^k\)，必须在 \(n\) 个因子中选出恰好 \(k\) 个，让它们贡献 \(y\)；剩下的 \(n-k\) 个因子贡献 \(x\)。从 \(n\) 个位置挑 \(k\) 个，有多少种挑法？\(\binom nk\) 种。于是

\[
(x+y)^n
=\sum_{k=0}^n\binom nk x^{n-k}y^k.
\]

这就是二项式定理的全部内容。系数 \(\binom nk\) 不是在算阶乘，而是在数选择——每项系数正好是``从 \(n\) 个因子中挑出 \(k\) 个贡献 \(y\)''的方案数。数选择这个动作天然等于组合数，组合数天然出现在多项式展开里。

你可能会想：用阶乘公式验证也可以，为什么偏要绕回``选择''？因为阶乘只告诉你等式成立，不告诉你为什么。选择视角让你看见新恒等式时能自己造出解释。

举个例子。在 \((x+y)^n\) 中令 \(x=y=1\)：

\[
\sum_{k=0}^n\binom nk=2^n.
\]

左边是把所有大小为 \(0,1,\ldots,n\) 的子集数一遍——你是在按子集大小分类，数所有子集。右边是让每个元素独立决定``进还是不进''——每个元素有 2 种选择，\(n\) 个元素就是 \(2^n\)。同一个对象（全体子集）被两种不同的分类方式数了，数目自然相等。代数上代人 \(x=y=1\) 只花了一行，但这个恒等式不是代数巧合，它是两种心智动作的对应：分类加总 vs 独立决策。

\subsection{Pascal 恒等式来自``某个元素选不选''}

有了这个视角，Pascal 恒等式也不再是三角里那个``肩上两数之和''的背诵口诀。

从 \(n\) 个人中选 \(k\) 个组队。随便挑一个人，叫他 A。所有的选法可以分成两类：

\begin{itemize}
\tightlist
\item
  不选 A：从剩下的 \(n-1\) 个人里选 \(k\) 个，有 \(\binom{n-1}k\) 种；
\item
  选 A：A 已经占了一个名额，从剩下的 \(n-1\) 个人里再选 \(k-1\) 个，有 \(\binom{n-1}{k-1}\) 种。
\end{itemize}

这两类互斥，且穷尽了所有可能（A 要么选上要么没选上，没有第三种状态）。于是

\[
\binom nk
=\binom{n-1}k+\binom{n-1}{k-1}.
\]

没有展开任何阶乘。你只是盯着集合里的一个固定元素，问了一句``选不选它''，恒等式就出来了。Pascal 三角里每一个数等于肩上两数之和，就是``固定一个元素，分类讨论''的直接翻译——递推和组合意义是同一件事的两个面。

\subsection{双重计数：同一集合数两遍}

上面的思路可以系统化为一种证明方法：双重计数。方法是先定义一个清楚的对象集合，然后按两种不同的特征各自数一遍。

回忆上一篇选委员会的例子——先从 \(n\) 人选 \(k\) 人当委员，再从委员中选一人当主席；另一种做法是先选主席，再从剩下的 \(n-1\) 人选 \(k-1\) 个委员。两种做法数的是同一批``带主席的委员会''，所以数目相等，导出了 \(k\binom nk=n\binom{n-1}{k-1}\)。这不是代数变形，而是两个视角在看同一堆东西。

再看一个例子。班上有 \(m\) 个男生和 \(n\) 个女生，要选 \(r\) 个人参加比赛。一种方法是不分男女，直接从 \(m+n\) 个人里选 \(r\) 个：

\[
\binom{m+n}r.
\]

另一种方法是按选中的男生数分类。选中的 \(r\) 个人里，可能有 0 个男生、1 个男生、……、最多 \(r\) 个男生。如果恰有 \(k\) 个男生，那么女生就要选 \(r-k\) 个。男生那边的选法是 \(\binom mk\)，女生那边是 \(\binom n{r-k}\)。把所有可能的 \(k\) 加总：

\[
\sum_k\binom mk\binom n{r-k}.
\]

两种方法数的是同一批人选，因此

\[
\binom{m+n}r
=\sum_k\binom mk\binom n{r-k}.
\]

这就是 Vandermonde 恒等式。注意 \(k\) 的实际范围是 \(0\le k\le r\) 且 \(k\le m\) 且 \(r-k\le n\)——超出这些范围的组合数自然为 0，不必在求和号上把边界写得很复杂。

拿个具体数字试一下：\(m=3\) 个男生，\(n=2\) 个女生，选 \(r=2\) 人。直接选：\(\binom 52=10\)。分类：0 个男生（即 2 个女生）\(\binom 30\binom 22=1\times1=1\)；1 个男生 1 个女生 \(\binom 31\binom 21=3\times2=6\)；2 个男生 0 个女生 \(\binom 32\binom 20=3\times1=3\)。合计 \(1+6+3=10\)，对上了。

\subsection{对称性也有对象解释}

恒等式

\[
\binom nk=\binom n{n-k}
\]

用阶乘验证只需要两行：\(n!/(k!(n-k)!)=n!/((n-k)!k!)\)。但你盯着它看，总觉得背后应该有一个更简单的理由。

确实有。从 \(n\) 个人里选 \(k\) 个去台上，等价于从 \(n\) 个人里选 \(n-k\) 个留在台下。每选出一个 \(k\) 人团队，就唯一对应一个 \(n-k\) 人的补集团队。这是一个双射——双方元素可以一一配对，不重不漏。所以两边数目必然相同。

这里的关键词是``补集''。你不必计算任何阶乘，只需要意识到``选谁上台''和``选谁不上台''数的是同一件事的两个侧面。组合恒等式证明的最高优先级是找双射或双重计数：如果能找到两个视角在看同一堆对象，恒等式就不需要代数验证；代数验证可以核对，却可能把藏在结构里的原因掩埋掉。

双射证明和代数证明的差别，就像你解释为什么镜子里的你和真实的你左右互换——一个人可能给你写一整页坐标系变换公式，另一个人让你站到镜子前指给你看。两者都没错，但后者让你觉得``我也可以想到''。
